\chapter{Лабораторная работа №1.1. Исследование функции одной переменной и визуализация функции двух переменных}
\label{ch:chap1}

\section{Цель работы}
\label{sec:goal1}

    Освоить методы численного анализа функций в MATLAB: вычисление производной и определённого интеграла, решение нелинейных уравнений, построение двумерных и трёхмерных графиков. Научиться применять встроенные функции \verb|diff|, \verb|integral|, \verb|fsolve|, \verb|meshgrid|, \verb|surf| для инженерных расчётов.

\section{Задание к лабораторной работе}
\label{sec:task1}

    В соответствии с методическими указаниями \cite{andreev2024tmo} и вариантом №2 необходимо:

    \begin{enumerate}
        \item Определить функцию $f(x)$, заданную по варианту №2:
        \begin{equation}
        \label{eq:f}
            f(x) = \frac{x^3}{\cos(x)}.
        \end{equation}
        \item Построить графики функции $f(x)$, её производной $f'(x)$ и интеграла $F(x) = \int_{x_0}^{x} f(y) \, dy$ на интервале $x \in [0.5; 4]$.
        \item Решить уравнение вида $a \cdot x + b = f(x)$ графическим и численным методами. Коэффициенты $a$ и $b$ задать самостоятельно. Использовать функцию \verb|fsolve| для численного решения.
        \item Определить функцию двух переменных $F(x,y)$ по варианту №2:
        \begin{equation}
        \label{eq:F}
            F(x,y) = \frac{\sqrt{x^2 + \dfrac{\pi}{y}}}{x^2}.
        \end{equation}
        \item Построить трёхмерный график $F(x,y)$ с использованием \verb|meshgrid| и \verb|surf|.
        \item Подписать все графики с помощью \verb|title|, \verb|xlabel|, \verb|ylabel|, добавить легенды \verb|legend|.
    \end{enumerate}

\section{Ход выполнения работы}
\label{sec:work1}

\subsection{Задание функции и построение графиков}
\label{subsec:fig1}

    Функция одной переменной задана выражением \eqref{eq:f}. Диапазон изменения аргумента: $x \in [0.5; 4]$ с шагом $0.05$. Следует отметить, что функция имеет разрыв в точке $x = \pi/2 \approx 1.5708$, где $\cos(x) = 0$. Для корректного построения графиков на этом участке шаг уменьшен, а разрыв обработан средствами MATLAB.

    Производная вычислена численно с помощью функции \verb|diff| как отношение приращений:
    \begin{equation}
    \label{eq:deriv}
        f'(x_i) \approx \frac{f(x_{i+1}) - f(x_i)}{x_{i+1} - x_i}.
    \end{equation}

    Интеграл вычислен с помощью функции \verb|integral| как определённый интеграл с переменным верхним пределом:
    \begin{equation}
    \label{eq:integ}
        F(x) = \int_{0.5}^{x} f(y) \, dy.
    \end{equation}

    \begin{figure}[ht]
        \centering
        \includegraphics[width=0.9\textwidth]{lab1_1_plot_function}
        \caption{Графики функции $f(x)$, её производной и интеграла}
        \label{fig:lab1_1_plot_function}
    \end{figure}

    На рисунке \ref{fig:lab1_1_plot_function} представлены три кривые: синяя — функция $f(x)$, красная — производная, зелёная — интеграл. Виден характерный разрыв функции в окрестности $x = \pi/2$.

\subsection{Решение уравнения $a \cdot x + b = f(x)$}
\label{subsec:eq1}

    Уравнение решено для $a = -2$, $b = 10$. Точный корень найден функцией \verb|fsolve| с начальной точкой $x_0 = 2$.

    \begin{figure}[ht]
        \centering
        \includegraphics[width=0.9\textwidth]{lab1_1_plot_equation}
        \caption{Графическое решение уравнения $a \cdot x + b = f(x)$}
        \label{fig:lab1_1_plot_equation}
    \end{figure}

    На рисунке \ref{fig:lab1_1_plot_equation} показаны: синяя кривая — функция $f(x)$, magenta — прямая $a \cdot x + b$, красная точка — найденный корень.

\subsection{Трёхмерный график функции двух переменных}
\label{subsec:fig2}

    Функция двух переменных задана выражением \eqref{eq:F}. Сетка аргументов построена в диапазоне $x, y \in [0.1; 5]$ с шагом $0.05$. Область определения ограничена условиями $y \neq 0$ и $x \neq 0$.

    \begin{figure}[ht]
        \centering
        \includegraphics[width=0.9\textwidth]{lab1_1_plot_3d}
        \caption{Трёхмерный график функции $F(x,y)$}
        \label{fig:lab1_1_plot_3d}
    \end{figure}

    На рисунке \ref{fig:lab1_1_plot_3d} представлена поверхность $F(x,y)$, построенная функцией \verb|surf|, с цветовой шкалой \verb|colorbar|.

\section{Выводы}
\label{sec:concl1}

    В ходе лабораторной работы №1.1 освоены методы численного анализа функций в MATLAB. Исследована функция $f(x) = x^3/\cos(x)$, имеющая разрыв второго рода в точке $x = \pi/2$. Вычислены её производная и интеграл, решено нелинейное уравнение $a \cdot x + b = f(x)$. Построена трёхмерная поверхность для функции двух переменных. Получены навыки работы с анонимными функциями, численными методами и визуализацией данных.

\endinput